% 6. Novelty of the Approach.
\section{Novelty of the Approach}

This project combines several individually known techniques---position-based
dynamics, GPU compute dispatch, and continuous level-of-detail crowd
simulation---into a unified real-time pipeline whose specific design choices
constitute the novel contribution.  This section identifies each distinct
novelty and explains why it is non-trivial.

\subsection{Density-Driven Physics LOD via the Spatial Hash}

The most significant novelty is the use of the spatial hash occupancy count as
a real-time density heatmap that drives physics fidelity allocation.  Prior
CLOD systems for crowd simulation universally determine LOD tier from the
agent's distance to a camera or focus point~\cite{ulicny2001,helbing1995}.
This approach is intuitive for rendering LOD---objects further away need fewer
polygons---but is poorly motivated for physics: a dense cluster of agents one
hundred metres from the camera still generates as many collision contacts as
the same cluster at ten metres, and resolving them correctly is equally
important for plausible crowd behaviour.  Conversely, a lone agent directly in
front of the camera in an open area generates no contacts regardless of solver
fidelity.

The density heatmap approach inverts this relationship.  Physics fidelity is
allocated where the solver has the most work to do: dense regions receive the
full $3\times 3$ neighbourhood search and up to twelve contact resolutions per
Jacobi pass, while sparse regions use only a centre-cell lookup with a single
contact.  This is both physically motivated and computationally efficient.

Critically, the density field is obtained at no extra cost.  The
\texttt{KernelHashAgents} kernel must atomically insert every agent into the
spatial hash regardless of any LOD system; the per-bucket counter that enables
this insertion is the density value.  No second pass, no additional buffer, and
no extra kernel are required.  A single additional read of
\texttt{\_HashGridCount[h]} inside \texttt{KernelAssignLOD} is sufficient to
assign the LOD tier.

\subsection{GPU-Native LOD Assignment}

In existing CPU-side implementations (including our own \textit{PBDCLODJobs}
variant), LOD tier assignment runs as a single-threaded \texttt{IJob} on the
CPU before the parallel constraint jobs are dispatched.  This introduces a
serial bottleneck that grows with agent count: the categorisation pass cannot
be parallelised over the CPU job system without expensive atomic writes into
shared index lists.

The compute shader implementation moves the entire LOD classification onto the
GPU.  \texttt{KernelAssignLOD} runs as a fully parallel dispatch---one thread
per agent---with no synchronisation other than the implicit barrier between
compute dispatches.  Each thread reads two integers (the agent's cell
coordinates and the bucket count) and writes one (the \texttt{lodLevel} field
into the agent buffer).  The kernel is latency-bound rather than
bandwidth-bound, giving the GPU scheduler maximum freedom to interleave it with
other work.

\subsection{Elimination of Frame-Offset Skipping}

Camera-distance CLOD systems commonly save computation for distant agents by
resolving their constraints only every $k$-th frame (in our \textit{PBDCLODJobs}
implementation, $k = 4$ for the \textsc{Low} tier).  While effective in terms
of raw throughput, this technique introduces visible temporal artefacts: agents
that skip constraint resolution accumulate penetration over the skipped frames,
producing a characteristic stuttering or ``popping'' motion that is especially
noticeable when those agents re-enter the camera's high-LOD zone.

The density-driven system eliminates frame-offset skipping entirely.  Every
agent resolves its constraints on every frame.  The computational savings come
instead from reducing the \emph{scope} of the constraint solve---search radius
and contact count---rather than its \emph{frequency}.  Because the spatial hash
is already built as part of the pipeline, the marginal cost of dispatching a
low-tier constraint thread is small: the thread exits after checking at most
one contact in a single bucket.  The result is smoother motion for all agents
at a modest increase in per-frame GPU occupancy.

\subsection{Single-Buffer Density Map with No Memory Overhead}

A predictive density map that is computed before LOD assignment and then reused
for neighbour lookup during constraint resolution would ordinarily require an
extra buffer and an extra kernel pass.  Our implementation avoids this by
recognising that the spatial hash count array already serves all three
purposes: atomic counter during hashing, density heatmap during LOD assignment,
and neighbourhood index during constraint resolution.  The
\texttt{\_HashGridCount} buffer is read in three separate kernels with no
intervening writes, so no copies or redundant allocations are needed.
% TODO: Quantify the memory saving relative to a naive two-buffer design
%       once final agent counts are fixed (N agents x 4 bytes per density entry).

\subsection{Asynchronous Readback with Version-Guarded Callbacks}

The \texttt{AsyncGPUReadback} API allows the GPU to continue rendering the
next frame while the previous frame's simulation results are transferred to the
CPU.  This decouples the physics timestep from the main-thread stall that would
result from a synchronous \texttt{GetData} call.  The one- to two-frame lag
in NavMesh position updates is imperceptible at interactive frame rates because
NavMesh pathfinding operates on a planning horizon of many frames.

The version counter guards against a subtle race condition: if the simulation
is respawned (e.g., because the user switched solver mode or changed the agent
count) between the \texttt{AsyncGPUReadback.Request} call and the callback
firing, the callback would apply stale agent data from the old simulation to
the new agent list.  The counter---incremented on every \texttt{Init}---is
captured at request time and compared inside the callback; a mismatch causes
the callback to return immediately without modifying any state.  To our
knowledge this guard pattern is not described in existing Unity compute shader
tutorials or documentation.